% ============================================================
% Chapter 5: Development
%   Section: News Acquisition: Discovery and Extraction
%     Subsection: Discovery Strategies
% Included from chapters/05-development/02-news-acquisition/news-acquisition.tex
% ============================================================

\subsection{Discovery Strategies}
\label{subsec:scraping-discovery}

Discovery tries three strategies per source, cheapest first, and stops
at the first that returns links (Table~\ref{tab:discovery-strategies}).
Every attempt is recorded under the purpose \texttt{discovery}, so a
feed that starts failing appears on the monitoring page next to the
articles it no longer delivers.

\begin{table}[ht]
\centering
\small
\caption{Discovery strategies, in the order they are tried.}
\label{tab:discovery-strategies}
\begin{tabularx}{\linewidth}{@{}lXl@{}}
\toprule
Strategy & Tried when & Gives per link \\
\midrule
1. RSS feed            & Always, if the source has a feed URL.
                       & time, title, summary \\
2. Feed discovery      & The feed gave nothing: it is searched for on
                         the feed URL and then on the homepage.
                       & URL only \\
3. Topic section pages & There is no feed at all.
                       & URL only \\
\bottomrule
\end{tabularx}
\end{table}

\paragraph{The RSS feed.} The feed is downloaded through the shared
fetcher, so it is guarded and has a timeout (the feed library's own
HTTP client has neither), and is then parsed with \texttt{feedparser}.
Each item keeps its publication time (in UTC), its title and a summary
reduced to 400 characters of plain text: what the candidate list shows
and the AI selection reads before any article is opened. On
4~October~2026, all 1,248 URLs from the 29 feeds that answered carried
a time.

\paragraph{Feed discovery.} When the configured feed yields nothing,
\texttt{trafilatura}'s feed discovery looks for one, first on the
configured feed URL and then on the homepage. It is lenient with
malformed XML and finds the feed a site advertises today rather than
the one it advertised when the source was configured.

\paragraph{Topic section pages.} Four configured feed URLs were found
to answer 404 (National Geographic, Reuters, RTVE and SINC), and their
homepages advertised no other feed, so those sources produced nothing
although their sites were up. Every news site files its articles under
sections, and the sections this project wants are the ones its topic
configuration already names. The third strategy
(\path{strategies/topic_pages.py}) reads those section pages as HTML
and keeps the article links on them:

\begin{enumerate}
    \item first, the sections the homepage itself links to (a short
    path, at most three segments, ending in a topic's section name), so
    that nothing is spent on a guess;
    \item then guesses, the base URL followed by a section name, one per
    topic, with at most six section pages per source in all, and none
    when the homepage already refused the request (EFE and Reuters
    refused every guess in the same way they refused the homepage).
\end{enumerate}

Spanish sources also get Spanish section names (\texttt{/ciencia/},
\texttt{/salud/}, \texttt{/cultura/}\ldots), kept in a separate table on
purpose, because the English table also decides whether a feed link
looks like an article, and that behaviour was already settled on the
English feeds. A link is kept when it is on the source's host or one
of its subdomains (NASA's science sections live on
\texttt{science.nasa.gov}), is not itself a section, and passes the
same article-shape test as feed links
(Section~\ref{subsec:scraping-filters}). A section page gives no title
or summary; the candidate list shows a title built from the link's
slug and marks it as such.

\paragraph{Stale feeds.} A feed whose items are all older than the
age limit counts as empty, so the next strategy runs. RTVE's feed has
carried only items from June~2022 since then; until this rule existed,
a Culture round listed two of them.

\paragraph{Concurrency.} Sources are discovered four at a time: feeds
are other people's servers and several are slow (two timed out in the
check of 23~September), and four at once keeps a round to seconds
without leaning on anyone. Results come back in the order the sources
were given, whatever order they finish in.

\paragraph{A regression found in use.} From 4~October, when discovery
started asking every strategy for publication times, the
topic-page strategy inherited the RSS strategy's feed reader and read
the dead feed again instead of crawling the sections: National
Geographic, Reuters and SINC found nothing for a day. It was fixed on
5~October, with a test that runs the strategy through
\texttt{DiscoveryService} as discovery does rather than in isolation.
